% GENERATED FILE. Edit canonical lessons, then run npm run generate:books.
% canonical-source-hash: fnv1a64:5fbf611893ace417
% canonical-lessons: MW-C117-kyu, MW-C117-karan, MW-C117-vastai, MW-C117-khatar, MW-C117-kyukai

\chapter{Why, Reason, For, Because of, Because}
\label{ch:mw-why-reason-for-because-of-because}
\hlchaptermodality{117}

\begin{chapteropening}
\textbf{By the end of this chapter:} \emph{I can say why, a reason, for, because of and because.}

Complete the last lesson of chapter 117: I can say why, a reason, for, because of and because.
\end{chapteropening}

\section[\texorpdfstring{\mw{क्यूं} (kyū̃)}{क्यूं (kyū̃)}]{\mw{क्यूं} (kyū̃) — why}

\label{lesson:MW-C117-kyu}



\begin{quote}
Before the new one: say the Marwadi for a right, then the Marwadi for a chance.
\end{quote}

\subsection*{You'll want to know: \mw{क्यूं}}
\textbf{\mw{क्यूं}} — \emph{kyū̃} — \textquotedblleft{}why\textquotedblright{}.

\begin{cousinweb}
\textbf{\mw{क्यूं}} (\emph{kyū̃}) is a tadbhava, a word inherited from Sanskrit and reshaped by sound change. It comes from \emph{kim}, \textquotedblleft{}what, why\textquotedblright{}. CDIAL lists Old Marwari \emph{kyu} under it (CDIAL 3164). Its \emph{k-} is the old question stem.
\end{cousinweb}

\begin{grammarlens}[title={What you've built}]
One more word for this chapter.
\end{grammarlens}

\subsection*{Your turn}
\begin{itemize}
\raggedright
  \item \emph{Say it:} \emph{kyū̃}
  \item \emph{Say it:} \emph{kyū̃}, once more
  \item \emph{Say it:} say \emph{kyū̃}
  \item \emph{Recall:} say the Marwadi for a right, then the Marwadi for a chance, then say \emph{kyū̃} again
\end{itemize}

\begin{tcolorbox}[breakable,colback=teal!4,colframe=teal!35!black,title={Before you move on}]
What does \mw{क्यूं} mean? (\textquotedblleft{}Why\textquotedblright{}.) Say it once more.
\end{tcolorbox}
\section[\texorpdfstring{\mw{कारण} (kāraṇ)}{कारण (kāraṇ)}]{\mw{कारण} (kāraṇ) — a reason}

\label{lesson:MW-C117-karan}



\begin{quote}
Before the new one: say the Marwadi for a chance, then the Marwadi for why.
\end{quote}

\subsection*{You'll want to know: \mw{कारण}}
\textbf{\mw{कारण}} — \emph{kāraṇ} — \textquotedblleft{}a reason\textquotedblright{}.

\begin{cousinweb}
\textbf{\mw{कारण}} (\emph{kāraṇ}) is Sanskrit \emph{kāraṇa}, \textquotedblleft{}cause\textquotedblright{}, taken over whole as a learned loan, a tatsama (CDIAL 3057). CDIAL also lists Gujarati \emph{kāraṇ} and Hindi \emph{kāran}, \textquotedblleft{}reason\textquotedblright{}, under it.
\end{cousinweb}

\begin{grammarlens}[title={What you've built}]
One more word for this chapter.
\end{grammarlens}

\subsection*{Your turn}
\begin{itemize}
\raggedright
  \item \emph{Say it:} \emph{kāraṇ}
  \item \emph{Say it:} \emph{kāraṇ}, once more
  \item \emph{Say it:} say \emph{kāraṇ}
  \item \emph{Recall:} say the Marwadi for a chance, then the Marwadi for why, then say \emph{kāraṇ} again
\end{itemize}

\begin{tcolorbox}[breakable,colback=teal!4,colframe=teal!35!black,title={Before you move on}]
What does \mw{कारण} mean? (\textquotedblleft{}A reason\textquotedblright{}.) Say it once more.
\end{tcolorbox}
\section[\texorpdfstring{\mw{वास्तै} (vāstai)}{वास्तै (vāstai)}]{\mw{वास्तै} (vāstai) — for, for the sake of}

\label{lesson:MW-C117-vastai}



\begin{quote}
Before the new one: say the Marwadi for why, then the Marwadi for a reason.
\end{quote}

\subsection*{You'll want to know: \mw{वास्तै}}
\textbf{\mw{वास्तै}} — \emph{vāstai} — \textquotedblleft{}for, for the sake of\textquotedblright{}.

\begin{cousinweb}
\textbf{\mw{वास्तै}} (\emph{vāstai}) is a loan from Arabic \emph{wāsiṭa}, \textquotedblleft{}means, intermediary\textquotedblright{}, which reached Marwadi through Persian. Its Arabic root is \emph{w-s-ṭ} (Lane). Hindi and Urdu \emph{vāste}, \textquotedblleft{}for the sake of\textquotedblright{}, grew from its oblique form; Marwadi ends it in \emph{-ai}.
\end{cousinweb}

\begin{grammarlens}[title={What you've built}]
One more word for this chapter.
\end{grammarlens}

\subsection*{Your turn}
\begin{itemize}
\raggedright
  \item \emph{Say it:} \emph{vāstai}
  \item \emph{Say it:} \emph{vāstai}, once more
  \item \emph{Say it:} say \emph{vāstai}
  \item \emph{Recall:} say the Marwadi for why, then the Marwadi for a reason, then say \emph{vāstai} again
\end{itemize}

\begin{tcolorbox}[breakable,colback=teal!4,colframe=teal!35!black,title={Before you move on}]
What does \mw{वास्तै} mean? (\textquotedblleft{}For, for the sake of\textquotedblright{}.) Say it once more.
\end{tcolorbox}
\section[\texorpdfstring{\mw{खातर} (khātar)}{खातर (khātar)}]{\mw{खातर} (khātar) — because of, for}

\label{lesson:MW-C117-khatar}



\begin{quote}
Before the new one: say the Marwadi for a reason, then the Marwadi for for.
\end{quote}

\subsection*{You'll want to know: \mw{खातर}}
\textbf{\mw{खातर}} — \emph{khātar} — \textquotedblleft{}because of, for\textquotedblright{}.

\begin{cousinweb}
\textbf{\mw{खातर}} (\emph{khātar}) is a loan from Arabic \emph{ḫāṭir}, \textquotedblleft{}mind; sake\textquotedblright{}, which reached Marwadi through Persian. Its Arabic root is \emph{ḫ-ṭ-r} (Lane). Persian uses it in \emph{ba xāṭir}, \textquotedblleft{}for the sake of\textquotedblright{}.
\end{cousinweb}

\begin{grammarlens}[title={What you've built}]
One more word for this chapter.
\end{grammarlens}

\subsection*{Your turn}
\begin{itemize}
\raggedright
  \item \emph{Say it:} \emph{khātar}
  \item \emph{Say it:} \emph{khātar}, once more
  \item \emph{Say it:} say \emph{khātar}
  \item \emph{Recall:} say the Marwadi for a reason, then the Marwadi for for, then say \emph{khātar} again
\end{itemize}

\begin{tcolorbox}[breakable,colback=teal!4,colframe=teal!35!black,title={Before you move on}]
What does \mw{खातर} mean? (\textquotedblleft{}Because of, for\textquotedblright{}.) Say it once more.
\end{tcolorbox}
\section[\texorpdfstring{\mw{क्यूंकै} (kyū̃kai)}{क्यूंकै (kyū̃kai)}]{\mw{क्यूंकै} (kyū̃kai) — because}

\label{lesson:MW-C117-kyukai}



\begin{quote}
Before the new one: say the Marwadi for for, then the Marwadi for because of.
\end{quote}

\subsection*{You'll want to know: \mw{क्यूंकै}}
\textbf{\mw{क्यूंकै}} — \emph{kyū̃kai} — \textquotedblleft{}because\textquotedblright{}.

\begin{cousinweb}
\textbf{\mw{क्यूंकै}} (\emph{kyū̃kai}) joins \emph{kyū̃}, \textquotedblleft{}why\textquotedblright{}, which CDIAL lists as Old Marwari \emph{kyu} under Sanskrit \emph{kim}, \textquotedblleft{}what\textquotedblright{} (CDIAL 3164), with \emph{kai}, \textquotedblleft{}that\textquotedblright{}. CDIAL takes the similar Hindi \emph{ki}, \textquotedblleft{}that\textquotedblright{}, from Persian (CDIAL 3164), but gives no Marwari form of it, so that part is not settled.
\end{cousinweb}

\begin{grammarlens}[title={What you've built}]
One more word for this chapter.
\end{grammarlens}

\subsection*{Your turn}
\begin{itemize}
\raggedright
  \item \emph{Say it:} \emph{kyū̃kai}
  \item \emph{Say it:} \emph{kyū̃kai}, once more
  \item \emph{Say it:} say \emph{kyū̃kai}
  \item \emph{Recall:} say the Marwadi for for, then the Marwadi for because of, then say \emph{kyū̃kai} again
\end{itemize}

\begin{tcolorbox}[breakable,colback=teal!4,colframe=teal!35!black,title={Before you move on}]
What does \mw{क्यूंकै} mean? (\textquotedblleft{}Because\textquotedblright{}.) Say it once more.
\end{tcolorbox}
